\begin{problem}[第30届全国中学生物理竞赛复赛][逆康普顿散射]{8}
光子被电子散射时，如果初态电子具有足够的动能，以至于在散射过程中有能量从电子转移到光子，则该散射被称为逆康普顿散射。当低能光子与高能电子发生对头碰撞时，就会出现逆康普顿散射。已知电子静止质量为 $m_e$，真空中的光速为 $c$。若能量为 $E_e$ 的电子与能量为 $E_\gamma$ 的光子相向对碰，
\begin{subproblem}
求散射后光子的能量；
\end{subproblem}
\begin{subproblem}
求逆康普顿散射能够发生的条件；
\end{subproblem}
\begin{subproblem}
如果入射光子能量为 $2.00\,\mathrm{eV}$，电子能量为 $1.00 \times 10^{9}\,\mathrm{eV}$，求散射后光子的能量。已知 $m_e = 0.511 \times 10^{6}\,\mathrm{eV}/c^2$。计算中有必要时可利用近似：如果 $|x| \ll 1$，有 $\sqrt{1-x} \approx 1 - \frac{1}{2}x$。
\end{subproblem}
\end{problem}